% Generated draft: requires mathematical and source review.
\documentclass[11pt]{article}
\usepackage[margin=25mm]{geometry}
\usepackage{fontspec}
\setmainfont{Latin Modern Roman}
\usepackage{amsmath,amssymb,mathtools,unicode-math}
\setmathfont{Latin Modern Math}
\usepackage{enumitem,needspace,xurl,hyperref}
\setlength{\parindent}{0pt}
\setlength{\parskip}{5pt}
\setlength{\emergencystretch}{3em}
\newcommand{\N}{\mathbb{N}}
\newcommand{\Z}{\mathbb{Z}}
\newcommand{\Q}{\mathbb{Q}}
\newcommand{\R}{\mathbb{R}}
\newcommand{\C}{\mathbb{C}}
\newcommand{\F}{\mathbb{F}}
\newcommand{\A}{\mathbb{A}}
\newcommand{\PP}{\mathbb{P}}
\newcommand{\E}{\mathbb{E}}
\newcommand{\eps}{\varepsilon}
\newcommand{\dd}{\,\mathrm{d}}
\begin{document}
\section*{UnsolvedMath ID: 30000780}
\textbf{Source problem code:} catalogue entry
\subsection{Definitions and mathematical statement}
Let $p\geq5$ be prime and let $E/F$ be an elliptic curve. Set \[F_\infty=F(E[p^\infty])=\bigcup_{r\geq1}F(E[p^r]),\qquad G=\operatorname{Gal}(F_\infty/F).\] Let $\rho:G\to\operatorname{GL}(V_\rho)$ be an irreducible self-dual Artin representation over $\overline{\Q}_p$, with $\dim V_\rho=d_\rho$. Here Artin means continuous with finite image, and self-dual means $V_\rho\simeq V_\rho^\vee$. The intended arithmetic interpretation requires $F$ to be a number field; this is not expressly stated in the extract. In that interpretation, choose a finite Galois extension $K/F$ inside $F_\infty$ through whose group $H=\operatorname{Gal}(K/F)$ the representation factors. With the usual $p$-primary Selmer group, set \[X(E/K)=\operatorname{Hom}_{\mathrm{cont}}(\operatorname{Sel}_{p^\infty}(E/K),\Q_p/\Z_p),\qquad s(E,\rho)=\dim_{\overline{\Q}_p}\operatorname{Hom}_H\bigl(V_\rho,X(E/K)\otimes_{\Z_p}\overline{\Q}_p\bigr).\] Thus the multiplicity is taken in the rationalized Pontryagin dual, not as an integral multiplicity in a torsion module, and it is not replaced by a Mordell--Weil multiplicity under an unstated finiteness assumption. Let $w(E,\rho)$ be the global root number of the Artin-twisted $L$-function of $E$, with its usual complex realization. The source's Artin-twist parity question is \[w(E,\rho)=(-1)^{s(E,\rho)}\ ?\] The clean/current text says a finite field through which $\rho$ factors. A representation factors through a finite quotient group, supplied here by a finite extension $K/F$; a literal finite-cardinality field is not the same object. This correction of interpretation is explicitly flagged rather than silently assumed. The supplied verification note attributes the restored identity to Conjecture 8 of the report. Its separate 2026-08-21 literature assessment says the passage only defines the multiplicity and root number and supplies no complete theorem resolving the implied parity question. Both qualifications are retained, together with the supplied open label.
\subsection{Short English statement}
For an irreducible self-dual finite-image representation in an elliptic curve's $p$-power division tower, ask whether the twisted root number equals the parity sign of its Selmer multiplicity. The extract's phrase finite field needs clarification as a finite extension giving the finite Galois quotient.
\subsection{Sources}
\begin{itemize}\item Local UnsolvedMath registry, permanent ID 30000780, source problem catalogue entry, \emph{Artin-Twist Parity for Elliptic Curves}, proposed year 2007; no individual proposer supplied. The verification note identifies Conjecture 8.\item Supplied citation: \emph{Algebraische Zahlentheorie} (2008), Oberwolfach Reports, DOI 10.4171/owr/2007/30, pp.\ 1741--1792: \url{https://doi.org/10.4171/owr/2007/30}. Supplied literature label: original Oberwolfach report.\end{itemize}
\end{document}
